\chapter{Capitolo 7 - Monoidali, enriched, fibrations, topos e dinamiche categoriali}

\section{1. Problema}

Con i Capitoli 4-6 abbiamo completato la transizione dai mattoni locali della teoria (oggetti, morfismi, composizione, identità) al livello trasformazionale (funtori, naturali, aggiunzioni, monadi/comonadi). Resta ora il passaggio verso le strutture categoriali avanzate che, nel corpus classico, consentono di modellare sistemi complessi multilivello:

\begin{enumerate}
\def\labelenumi{\arabic{enumi}.}
\tightlist
\item
  categorie monoidali;
\item
  categorie arricchite;
\item
  fibrations;
\item
  topos.
\end{enumerate}

Nel quadro classico queste strutture estendono potenza espressiva e generalità. In OCT dobbiamo aggiungere una condizione ulteriore: la complessità strutturale non deve diventare opacità ordinativa. Una struttura più ricca è utile solo se:

\begin{enumerate}
\def\labelenumi{\arabic{enumi}.}
\tightlist
\item
  resta interpretabile in termini di \texttt{Coh/Phi/Delta};
\item
  mantiene robustezza sotto variazioni contestuali;
\item
  non nasconde degenerazioni dietro simmetrie formali.
\end{enumerate}

Il problema del capitolo è quindi:

\textbf{come estendere le strutture avanzate della category theory mantenendo insieme rigore classico, selettività ordinativa e tracciabilità operativa?}

La difficoltà è concreta:

\begin{enumerate}
\def\labelenumi{\arabic{enumi}.}
\tightlist
\item
  strutture monoidali possono imporre simmetrie non sostenibili in dominio reale;
\item
  arricchimenti possono aumentare espressività ma anche amplificare rumore;
\item
  fibrations possono modellare dipendenze verticali ma introdurre perdita tra base e fibra;
\item
  topos possono offrire quadro logico interno potente ma non garantire, da soli, validità ordinativa.
\end{enumerate}

Questo capitolo fornisce quindi una grammatica avanzata, ma disciplinata da un principio unico:

\begin{enumerate}
\def\labelenumi{\arabic{enumi}.}
\tightlist
\item
  ogni estensione strutturale deve dichiarare il proprio costo ordinativo.
\end{enumerate}

\subsection{\texorpdfstring{1.1 Perché questo capitolo chiude la milestone \texttt{v1.3}}{1.1 Perché questo capitolo chiude la milestone v1.3}}

La milestone \texttt{v1.3} mira a espandere sistematicamente il corpus classico in chiave OCT (Capitoli 4-7). Il Capitolo 7 è il punto di chiusura perché affronta i costrutti dove la complessità teorica diventa maggiore e dove il rischio di ``equivalenza apparente'' cresce.

Se il capitolo regge:

\begin{enumerate}
\def\labelenumi{\arabic{enumi}.}
\tightlist
\item
  il ponte verso protocollo, benchmark e governance (\texttt{v1.4}) diventa naturale;
\item
  la teoria può passare da fondazione/estensione a standard operativo.
\end{enumerate}

\subsection{1.2 Obiettivo tecnico}

Obiettivo minimo del capitolo:

\begin{enumerate}
\def\labelenumi{\arabic{enumi}.}
\tightlist
\item
  definire criteri ordinativi per monoidali, enriched, fibrations, topos;
\item
  mostrare casi tipici in cui la struttura classica resta valida ma l'uso ordinativo richiede restrizioni;
\item
  fissare regole di accettazione per strutture avanzate in scenari reali.
\end{enumerate}

\begin{center}\rule{0.5\linewidth}{0.5pt}\end{center}

\section{2. Tesi del capitolo}

La tesi è articolata in dieci enunciati.

\begin{enumerate}
\def\labelenumi{\arabic{enumi}.}
\tightlist
\item
  La struttura monoidale classica è necessaria ma non sempre sufficiente in contesti ordinativi asimmetrici.
\item
  Le categorie arricchite diventano ordinativamente utili solo se l'arricchimento non distrugge margine coerentivo.
\item
  Le fibrations sono strumenti potenti per stratificare contesti, ma richiedono controllo esplicito del trasporto verticale.
\item
  Un topos può fornire logica interna forte senza garantire automaticamente realtà ordinativa dei costrutti.
\item
  Esiste una nozione di ``coerenza multilivello'' che collega base, fibre e dinamica globale.
\item
  Le simmetrie e dualità avanzate restano condizionate dal dominio osservativo.
\item
  La complessità strutturale deve essere accompagnata da criteri di fallimento tracciabili.
\item
  Le strutture avanzate possono essere classificate con la stessa grammatica A/B/C se si esplicitano i mapping.
\item
  L'estensione proposta resta compatibile con recupero classico.
\item
  Il capitolo prepara direttamente il passaggio metodologico dei Capitoli 12-14.
\end{enumerate}

\begin{center}\rule{0.5\linewidth}{0.5pt}\end{center}

\section{3. Definizioni canoniche}

\subsection{Definizione 7.1 (Categoria monoidale ordinativa)}

Una categoria monoidale ordinativa è una categoria monoidale classica \texttt{(C,\ tensor,\ I)} con vincolo aggiuntivo:

\begin{enumerate}
\def\labelenumi{\arabic{enumi}.}
\tightlist
\item
  le operazioni \texttt{tensor} su oggetti/morfismi rilevanti non devono produrre caduta sistematica sotto soglia su \texttt{Coh};
\item
  la composizione monoidale deve preservare \texttt{Phi} non nulla su classi operative dichiarate.
\end{enumerate}

\subsection{Definizione 7.2 (Simmetria monoidale condizionata)}

La simmetria \texttt{sigma\_\{A,B\}:\ A\ tensor\ B\ -\textgreater{}\ B\ tensor\ A} è ordinativamente ammissibile solo se il cambio di ordine non altera in modo critico invarianti ordinativi del dominio.

\subsection{Definizione 7.3 (Categoria arricchita ordinativa)}

Una categoria arricchita ordinativa è una categoria arricchita classica in cui l'oggetto-hom arricchito incorpora anche indicatori ordinativi minimi (coerenza, emergenza, margine).

\subsection{Definizione 7.4 (Fibration ordinativa)}

Una fibration ordinativa \texttt{p:\ E-\textgreater{}B} è una fibration classica in cui:

\begin{enumerate}
\def\labelenumi{\arabic{enumi}.}
\tightlist
\item
  il sollevamento cartesiano non introduce perdita ordinativa non dichiarata;
\item
  la relazione base-fibra mantiene coerenza multilivello su classi test.
\end{enumerate}

\subsection{Definizione 7.5 (Coerenza multilivello)}

Dato un oggetto fibrato, la coerenza multilivello è la compatibilità tra:

\begin{enumerate}
\def\labelenumi{\arabic{enumi}.}
\tightlist
\item
  coerenza nel livello base;
\item
  coerenza nel livello fibra;
\item
  coerenza delle mappe di passaggio.
\end{enumerate}

\subsection{Definizione 7.6 (Topos ordinativo)}

Un topos ordinativo è un topos classico usato in contesto OCT con vincolo che i costrutti logici interni impiegati in inferenze operative siano validati anche in termini ordinativi.

\subsection{Definizione 7.7 (Dinamica categoriale ordinativa)}

Una dinamica categoriale ordinativa è una evoluzione temporale di strutture/funtori in cui:

\begin{enumerate}
\def\labelenumi{\arabic{enumi}.}
\tightlist
\item
  i passaggi temporali restano formalmente leciti;
\item
  l'evoluzione non produce deriva sistematica verso sterilità/degenerazione.
\end{enumerate}

\subsection{Definizione 7.8 (Costo di complessità strutturale)}

Definiamo costo sintetico di complessità:

\texttt{Cost\_str\ :=\ c1*Instab\ +\ c2*Loss\_tr\ +\ c3*Opac}

dove \texttt{Opac} misura perdita di leggibilità causale del sistema al crescere della stratificazione.

\begin{center}\rule{0.5\linewidth}{0.5pt}\end{center}

\section{4. Sviluppo formale}

\subsection{4.1 Monoidali: potenza compositiva e asimmetria reale}

La struttura monoidale classica permette di modellare composizioni parallele o tensoriali. In OCT questa potenza è preziosa, ma va gestita con prudenza:

\begin{enumerate}
\def\labelenumi{\arabic{enumi}.}
\tightlist
\item
  in domini simmetrici il tensor può preservare bene \texttt{Coh/Phi};
\item
  in domini direzionali la simmetrizzazione può distruggere informazione funzionale.
\end{enumerate}

Da qui una regola operativa:

\begin{enumerate}
\def\labelenumi{\arabic{enumi}.}
\tightlist
\item
  \texttt{tensor} è ammesso sempre formalmente;
\item
  è ammesso in senso ordinativo solo dopo test di simmetria condizionata.
\end{enumerate}

Errore evitato:

\begin{enumerate}
\def\labelenumi{\arabic{enumi}.}
\tightlist
\item
  imporre simmetria per eleganza algebrica quando il dominio è strutturalmente orientato.
\end{enumerate}

\subsection{4.2 Arricchimento: espressività contro opacità}

Arricchire una categoria significa dare struttura più ricca agli hom. In OCT questo può essere vantaggioso perché:

\begin{enumerate}
\def\labelenumi{\arabic{enumi}.}
\tightlist
\item
  rende espliciti gradi di qualità relazionale;
\item
  consente metriche più fini su trasporto e composizione.
\end{enumerate}

Ma c'è un rischio:

\begin{enumerate}
\def\labelenumi{\arabic{enumi}.}
\tightlist
\item
  over-arricchimento che aumenta parametri e riduce stabilità inferenziale.
\end{enumerate}

Per questo proponiamo un criterio di arricchimento sobrio:

\begin{enumerate}
\def\labelenumi{\arabic{enumi}.}
\tightlist
\item
  introdurre solo componenti che migliorano decisione A/B/C;
\item
  evitare componenti decorative senza impatto diagnostico.
\end{enumerate}

\subsection{4.3 Fibrations: verticalità controllata}

Le fibrations permettono di separare livelli di descrizione (base/fibra). In OCT sono utili per modellare contesti multilivello, ma possono introdurre perdita verticale:

\begin{enumerate}
\def\labelenumi{\arabic{enumi}.}
\tightlist
\item
  informazioni mantenute in fibra ma non nella base;
\item
  coerenza locale alta in fibra con incoerenza globale in base;
\item
  trasporto verticale formalmente corretto ma ordinativamente dissipativo.
\end{enumerate}

Serve quindi un controllo bilanciato:

\begin{enumerate}
\def\labelenumi{\arabic{enumi}.}
\tightlist
\item
  test di coerenza interna alle fibre;
\item
  test di coerenza tra fibre e base;
\item
  test di stabilità dei lift sotto perturbazione contestuale.
\end{enumerate}

\subsection{4.4 Topos: logica interna e validità ordinativa}

Un topos fornisce ambiente logico potente (subobject classifier, exponentials, ecc.). In OCT questo non viene negato. Tuttavia:

\begin{enumerate}
\def\labelenumi{\arabic{enumi}.}
\tightlist
\item
  validità logica interna non equivale automaticamente a validità ordinativa esterna;
\item
  inferenze interne devono essere collegate a criteri di \texttt{Coh/Phi} quando usate per decisioni operative.
\end{enumerate}

Ne segue una distinzione:

\begin{enumerate}
\def\labelenumi{\arabic{enumi}.}
\tightlist
\item
  verità interna al topos;
\item
  accettabilità ordinativa nel dominio osservato.
\end{enumerate}

Questa distinzione evita che la ricchezza logica mascheri fragilità dinamiche.

\subsection{4.5 Dinamiche categoriali: dal statico al processo}

Molte strutture avanzate sono presentate staticamente. In applicazioni reali, però, categorie/funtori cambiano nel tempo. OCT richiede quindi una nozione dinamica:

\begin{enumerate}
\def\labelenumi{\arabic{enumi}.}
\tightlist
\item
  sequenze di strutture \texttt{C\_t};
\item
  transizioni \texttt{T\_t:\ C\_t\ -\textgreater{}\ C\_\{t+1\}};
\item
  monitoraggio di deriva su invarianti globali.
\end{enumerate}

Punto chiave:

\begin{enumerate}
\def\labelenumi{\arabic{enumi}.}
\tightlist
\item
  una struttura può essere ordinativamente valida a \texttt{t0} e non esserlo a \texttt{t1};
\item
  la validazione deve quindi includere traiettoria, non solo istanti.
\end{enumerate}

\subsection{4.6 Interazione tra monoidali e arricchimenti}

Quando si combinano tensor e arricchimento aumentano le interazioni non lineari:

\begin{enumerate}
\def\labelenumi{\arabic{enumi}.}
\tightlist
\item
  il tensor può amplificare rumore dei valori arricchiti;
\item
  un arricchimento utile localmente può diventare instabile in composizione tensoriale.
\end{enumerate}

Serve un test combinato:

\begin{enumerate}
\def\labelenumi{\arabic{enumi}.}
\tightlist
\item
  validità di \texttt{tensor} su oggetti semplici;
\item
  validità su oggetti arricchiti;
\item
  confronto differenziale tra i due regimi.
\end{enumerate}

\subsection{4.7 Interazione tra fibrations e topos}

In sistemi complessi è frequente usare logica interna (stile topos) su strutture stratificate (stile fibration). Qui il rischio maggiore è la disconnessione:

\begin{enumerate}
\def\labelenumi{\arabic{enumi}.}
\tightlist
\item
  coerenza logica locale alta;
\item
  incoerenza di trasporto tra livelli.
\end{enumerate}

Per evitarla proponiamo il principio di allineamento verticale:

\begin{enumerate}
\def\labelenumi{\arabic{enumi}.}
\tightlist
\item
  ogni inferenza interna usata operativamente deve avere tracciamento su livello base/fibra.
\end{enumerate}

\subsection{4.8 Criterio di accettazione per strutture avanzate}

Una struttura avanzata (\texttt{monoidal}, \texttt{enriched}, \texttt{fibration}, \texttt{topos}) è ordinativamente accettabile se:

\begin{enumerate}
\def\labelenumi{\arabic{enumi}.}
\tightlist
\item
  definita e verificata classicamente;
\item
  testata su invarianti ordinativi minimi;
\item
  robusta su variazioni contestuali rilevanti;
\item
  dotata di mappa fallimenti specifica.
\end{enumerate}

Questa regola unifica i quattro blocchi in una policy comune.

\subsection{4.9 Mappa fallimenti specifica del capitolo}

\begin{enumerate}
\def\labelenumi{\arabic{enumi}.}
\tightlist
\item
  \textbf{F-Mon2}: tensor formalmente corretto con simmetria ordinativamente distruttiva.
\item
  \textbf{F-Enr1}: arricchimento informativo ma instabile su composizione.
\item
  \textbf{F-Fib1}: lift cartesiano con perdita verticale critica.
\item
  \textbf{F-Fib2}: coerenza di fibra alta, coerenza di base bassa.
\item
  \textbf{F-Top1}: inferenza interna valida, ma non trasferibile operativamente.
\item
  \textbf{F-Dyn1}: validità puntuale alta, deriva temporale verso classe C.
\end{enumerate}

Questa tipologia è utile per collegare teoria avanzata e benchmark dedicati.

\subsection{4.10 Controllo di complessità: quando fermarsi}

Una lezione pratica del capitolo è che più struttura non significa sempre più scienza. Se \texttt{Cost\_str} cresce senza migliorare decisione o robustezza, l'estensione va ridotta.

Regola pragmatica:

\begin{enumerate}
\def\labelenumi{\arabic{enumi}.}
\tightlist
\item
  aggiungere complessità solo se migliora almeno uno tra:

  \begin{itemize}
  \tightlist
  \item
    precisione diagnostica;
  \item
    robustezza;
  \item
    trasferibilità.
  \end{itemize}
\end{enumerate}

Se nessun miglioramento è misurabile, la complessità è epistemicamente ingiustificata.

\subsection{4.11 Compatibilità con recupero classico}

Anche qui vale F10:

\begin{enumerate}
\def\labelenumi{\arabic{enumi}.}
\tightlist
\item
  disattivando il layer ordinativo, restano monoidali/enriched/fibrations/topos classici;
\item
  attivando il layer, si ottiene filtro di realtà compositiva senza negazione del quadro classico.
\end{enumerate}

\subsection{4.12 Condizione di non-ridondanza del capitolo}

Il capitolo è non ridondante se esiste almeno un caso in cui:

\begin{enumerate}
\def\labelenumi{\arabic{enumi}.}
\tightlist
\item
  la struttura avanzata classica è valida;
\item
  la stessa struttura fallisce in senso ordinativo per motivi tracciabili.
\end{enumerate}

Questa condizione è soddisfatta dai pattern F-Mon2/F-Fib1/F-Top1/D04/D05 e dai casi applicativi ad alta stratificazione.

\subsection{4.13 Matrice stratificata di validazione avanzata}

Per mantenere confrontabili analisi diverse, proponiamo una matrice di validazione a tre strati:

\begin{enumerate}
\def\labelenumi{\arabic{enumi}.}
\tightlist
\item
  \textbf{Strato Formale (SF)}: verifica classica della struttura avanzata;
\item
  \textbf{Strato Ordinativo (SO)}: verifica su \texttt{Coh/Phi/Delta}, margine e robustezza;
\item
  \textbf{Strato Dinamico (SD)}: verifica della stabilità temporale su evoluzioni non banali.
\end{enumerate}

Decisione proposta:

\begin{enumerate}
\def\labelenumi{\arabic{enumi}.}
\tightlist
\item
  \texttt{accepted\_strong} solo se SF, SO e SD sono positivi;
\item
  \texttt{accepted\_limited} se SF e SO positivi ma SD incompleto (uso limitato e monitorato);
\item
  \texttt{revise} se SF positivo ma SO o SD negativi;
\item
  \texttt{reject} se SF negativo o fallimenti ordinativi non recuperabili.
\end{enumerate}

Questa matrice evita un problema frequente: strutture avanzate formalmente corrette promosse troppo presto senza evidenza dinamica.

\subsection{4.14 Coupling tra tensor e arricchimento}

In molte applicazioni, monoidalità e arricchimento vengono usati insieme. Il coupling tra i due livelli può produrre effetti non lineari.

Proponiamo di monitorare:

\texttt{Kappa\_enr\_tensor\ :=\ Sens\_tensor(Enr)\ +\ Sens\_enr(tensor)}.

Interpretazione:

\begin{enumerate}
\def\labelenumi{\arabic{enumi}.}
\tightlist
\item
  \texttt{Kappa} basso: interazione controllata;
\item
  \texttt{Kappa} medio: interazione utile ma sensibile;
\item
  \texttt{Kappa} alto: rischio di instabilità con piccole variazioni di configurazione.
\end{enumerate}

Azioni consigliate quando \texttt{Kappa} è alto:

\begin{enumerate}
\def\labelenumi{\arabic{enumi}.}
\tightlist
\item
  ridurre dimensionalità dell'arricchimento;
\item
  restringere classe di oggetti su cui applicare \texttt{tensor};
\item
  introdurre normalizzazione intermedia.
\end{enumerate}

Questo blocco migliora la diagnosi preventiva: invece di reagire al collasso, si monitorano indicatori che lo anticipano.

\subsection{4.15 Regola di trasporto base-fibra nei sistemi fibrati}

Per le fibrations in contesti reali proponiamo una regola di trasporto in due fasi.

\begin{enumerate}
\def\labelenumi{\arabic{enumi}.}
\tightlist
\item
  \textbf{Fase fibra}: validazione locale dei costrutti e delle inferenze;
\item
  \textbf{Fase base}: validazione dell'impatto aggregato sul livello superiore.
\end{enumerate}

Solo il passaggio congiunto delle due fasi autorizza inferenze operative forti.\\
Questo impedisce due errori classici:

\begin{enumerate}
\def\labelenumi{\arabic{enumi}.}
\tightlist
\item
  assumere che ``vero in fibra'' implichi ``robusto in base'';
\item
  imporre vincoli di base che annullano informazione utile delle fibre.
\end{enumerate}

La regola rende esplicita la responsabilità di ogni livello nella tenuta complessiva.

\subsection{4.16 Protocollo minimo avanzato (PMA-7)}

Per standardizzare la verifica di strutture avanzate proponiamo un protocollo in sette passi:

\begin{enumerate}
\def\labelenumi{\arabic{enumi}.}
\tightlist
\item
  definire struttura classica e dominio d'uso;
\item
  dichiarare \texttt{Omega}, soglie e invarianti critici;
\item
  verificare SF (assiomi classici della struttura);
\item
  verificare SO (metriche ordinative);
\item
  verificare SD (stabilità temporale);
\item
  classificare fallimenti secondo F-Mon/F-Enr/F-Fib/F-Top/F-Dyn;
\item
  assegnare stato finale con motivazione e limiti dichiarati.
\end{enumerate}

Il PMA-7 non sostituisce prove matematiche complete, ma uniforma la fase di validazione comparativa e prepara direttamente i capitoli metodologici successivi.

\begin{center}\rule{0.5\linewidth}{0.5pt}\end{center}

\section{5. Proposizioni e teoremi locali}

\subsection{Proposizione 7.1 (Simmetria monoidale condizionata)}

Enunciato:
La simmetria monoidale è ordinativamente ammissibile solo su domini che non perdono informazione funzionale sotto scambio.

Intuizione:
domini direzionali possono violare stabilità pur mantenendo validità classica.

Stato: \texttt{in\_review}.

\subsection{Proposizione 7.2 (Necessità di coerenza multilivello nelle fibrations)}

Enunciato:
La sola correttezza del lift cartesiano non garantisce affidabilità ordinativa della stratificazione.

Intuizione:
serve compatibilità simultanea base-fibra-passaggio.

Stato: \texttt{validated}.

\subsection{Teorema 7.3 (Vincolo di allineamento verticale)}

Ipotesi:

\begin{enumerate}
\def\labelenumi{\arabic{enumi}.}
\tightlist
\item
  sistema fibrato con inferenze interne operative;
\item
  tracciamento completo base/fibra;
\item
  assenza di perdita verticale critica.
\end{enumerate}

Tesi:
Le inferenze stratificate sono ordinativamente affidabili nel dominio dichiarato.

Proof sketch:

\begin{enumerate}
\def\labelenumi{\arabic{enumi}.}
\tightlist
\item
  il tracciamento elimina ambiguità di livello;
\item
  l'assenza di perdita critica preserva coerenza multilivello;
\item
  segue affidabilità operativa dell'inferenza.
\end{enumerate}

Conseguenze:
criterio pratico per usare fibrations in pipeline reali.

Stato: \texttt{in\_proof}.

\subsection{Teorema 7.4 (Deriva dinamica di strutture avanzate)}

Ipotesi:

\begin{enumerate}
\def\labelenumi{\arabic{enumi}.}
\tightlist
\item
  sequenza temporale \texttt{C\_t} con transizioni \texttt{T\_t};
\item
  monitoraggio su invarianti globali;
\item
  presenza di accumulo negativo oltre soglia.
\end{enumerate}

Tesi:
Una struttura avanzata può passare da ordinativamente valida a non valida pur restando formalmente lecita.

Proof sketch:

\begin{enumerate}
\def\labelenumi{\arabic{enumi}.}
\tightlist
\item
  validità classica è puntuale e sintattica;
\item
  accumulo dinamico può degradare margine ordinativo;
\item
  la dinamica produce crossing verso classe B/C.
\end{enumerate}

Conseguenze:
necessità di validazione temporale, non solo statica.

Stato: \texttt{in\_proof}.

\subsection{Proposizione 7.5 (Topos non implica realtà ordinativa)}

Enunciato:
La ricchezza logica interna di un topos non implica automaticamente robustezza ordinativa dei costrutti usati.

Intuizione:
verità interna e operatività esterna sono livelli distinti.

Stato: \texttt{validated}.

\subsection{Teorema 7.6 (Criterio minimo di accettazione per strutture avanzate)}

Ipotesi:

\begin{enumerate}
\def\labelenumi{\arabic{enumi}.}
\tightlist
\item
  validità classica della struttura;
\item
  test ordinativo su invarianti;
\item
  robustezza contestuale;
\item
  mappa fallimenti specifica disponibile.
\end{enumerate}

Tesi:
La struttura può essere accettata per uso operativo ordinativo.

Proof sketch:

\begin{enumerate}
\def\labelenumi{\arabic{enumi}.}
\tightlist
\item
  requisito classico garantisce correttezza formale;
\item
  requisiti ordinativi garantiscono realtà compositiva;
\item
  la mappa fallimenti garantisce falsificabilità.
\end{enumerate}

Conseguenze:
fornisce standard unico per monoidali/enriched/fibrations/topos.

Stato: \texttt{in\_proof}.

\subsection{Proposizione 7.7 (Utilità della matrice SF/SO/SD)}

Enunciato:
La matrice stratificata SF/SO/SD riduce errori di attribuzione causale nei fallimenti di strutture avanzate.

Intuizione:
separando forma, validità ordinativa e dinamica, si evita di confondere errori assiomatici con instabilità applicative.

Stato: \texttt{in\_review}.

\subsection{Teorema 7.8 (Necessità del controllo dinamico per accettazione forte)}

Ipotesi:

\begin{enumerate}
\def\labelenumi{\arabic{enumi}.}
\tightlist
\item
  struttura avanzata con SF e SO positivi;
\item
  dominio con evoluzione temporale rilevante;
\item
  assenza di verifica SD.
\end{enumerate}

Tesi:
Non è possibile dichiarare \texttt{accepted\_strong} senza controllo dinamico.

Proof sketch:

\begin{enumerate}
\def\labelenumi{\arabic{enumi}.}
\tightlist
\item
  SF e SO descrivono correttezza e affidabilità locale;
\item
  in domini dinamici può emergere deriva non osservabile staticamente;
\item
  quindi SD è condizione necessaria per accettazione forte.
\end{enumerate}

Conseguenze:
formalizza il passaggio da validazione statica a validazione evolutiva.

Stato: \texttt{in\_proof}.

\subsection{Proposizione 7.9 (Monitoraggio del coupling come prevenzione)}

Enunciato:
Il monitoraggio di \texttt{Kappa\_enr\_tensor} riduce il rischio di instabilità non lineari in strutture monoidali arricchite.

Intuizione:
la sensibilità incrociata alta è spesso un segnale precoce di collasso su catene compositive complesse.

Stato: \texttt{in\_review}.

\begin{center}\rule{0.5\linewidth}{0.5pt}\end{center}

\section{6. Implicazioni operative}

\subsection{6.1 Per progettazione di architetture complesse}

Quando si introducono strutture avanzate, è utile decidere in anticipo:

\begin{enumerate}
\def\labelenumi{\arabic{enumi}.}
\tightlist
\item
  quale guadagno diagnostico ci si aspetta;
\item
  quale costo strutturale è accettabile;
\item
  quale criterio farà scattare revisione o rollback.
\end{enumerate}

\subsection{6.2 Per benchmark dedicati}

È consigliato preparare suite specifiche:

\begin{enumerate}
\def\labelenumi{\arabic{enumi}.}
\tightlist
\item
  \texttt{MON\_SUITE} per test tensor/simmetria;
\item
  \texttt{ENR\_SUITE} per stabilità arricchita;
\item
  \texttt{FIB\_SUITE} per coerenza multilivello;
\item
  \texttt{TOP\_SUITE} per trasferibilità inferenze interne;
\item
  \texttt{DYN\_SUITE} per deriva temporale.
\end{enumerate}

\subsection{6.3 Per reporting scientifico}

Ogni risultato su strutture avanzate dovrebbe riportare:

\begin{enumerate}
\def\labelenumi{\arabic{enumi}.}
\tightlist
\item
  definizione classica usata;
\item
  vincolo ordinativo applicato;
\item
  classe di fallimento osservata (se presente);
\item
  stato finale con motivazione.
\end{enumerate}

\subsection{\texorpdfstring{6.4 Per continuità con \texttt{v1.4}}{6.4 Per continuità con v1.4}}

Il capitolo fornisce già i mattoni metodologici dei Capitoli 12-14:

\begin{enumerate}
\def\labelenumi{\arabic{enumi}.}
\tightlist
\item
  metrica multilivello;
\item
  design benchmark orientato a fallimenti reali;
\item
  governance del costo di complessità.
\end{enumerate}

\subsection{6.5 Per uso interdisciplinare}

In contesti non puramente matematici, questa parte aiuta a evitare due estremi:

\begin{enumerate}
\def\labelenumi{\arabic{enumi}.}
\tightlist
\item
  riduzionismo formale (solo eleganza strutturale);
\item
  empirismo scollegato (solo risultato numerico senza struttura).
\end{enumerate}

La combinazione di entrambi è il contributo operativo di OCT in strutture avanzate.

\subsection{6.6 Per manutenzione del manoscritto}

Ogni modifica a questo capitolo dovrebbe propagare controlli almeno su:

\begin{enumerate}
\def\labelenumi{\arabic{enumi}.}
\tightlist
\item
  claim ledger dei Capitoli 6-7;
\item
  protocollo di validazione del Capitolo 12;
\item
  checklist benchmark del Capitolo 13.
\end{enumerate}

Questo evita disallineamenti tra teoria avanzata e infrastruttura metodologica.

\subsection{6.7 Per policy di rilascio progressivo}

In progetti applicativi è utile adottare una policy a maturità progressiva:

\begin{enumerate}
\def\labelenumi{\arabic{enumi}.}
\tightlist
\item
  fase esplorativa: strutture con stato \texttt{in\_review} ammesse solo in sandbox;
\item
  fase pre-rilascio: richiesti SF+SO positivi con limiti espliciti;
\item
  fase critica: richiesti SF+SO+SD con margine ordinativo minimo garantito.
\end{enumerate}

Questa policy riduce il rischio di rilasciare strutture eleganti ma non robuste.

\subsection{6.8 Per documentazione e peer review}

Per aumentare chiarezza è consigliabile allegare, per ogni struttura avanzata:

\begin{enumerate}
\def\labelenumi{\arabic{enumi}.}
\tightlist
\item
  tabella SF/SO/SD;
\item
  classe di fallimento prevalente;
\item
  stato finale (\texttt{accepted\_limited}, \texttt{accepted\_strong}, \texttt{revise}, \texttt{reject});
\item
  azione correttiva proposta.
\end{enumerate}

Questo formato aiuta revisori con background diversi a convergere su valutazioni comparabili.
Inoltre rende più trasparente il rapporto tra eleganza formale e robustezza operativa.

\begin{center}\rule{0.5\linewidth}{0.5pt}\end{center}

\section{7. Limiti e non-claim}

Limiti espliciti:

\begin{enumerate}
\def\labelenumi{\arabic{enumi}.}
\tightlist
\item
  il capitolo non contiene prove complete chiuse per tutte le varianti avanzate;
\item
  la metrica \texttt{Cost\_str} è proposta in forma minima, da specializzare per dominio;
\item
  la trattazione dei topos ordinativi è intenzionalmente prudente e non esaustiva.
\end{enumerate}

Failure mode riconosciuti:

\begin{enumerate}
\def\labelenumi{\arabic{enumi}.}
\tightlist
\item
  introdurre arricchimento senza guadagno misurabile;
\item
  usare fibrations senza controllo base-fibra;
\item
  assumere simmetria monoidale in domini direzionali;
\item
  validare strutture statiche ignorando deriva temporale;
\item
  usare logica interna del topos senza test di trasferibilità operativa.
\end{enumerate}

Box C - Non-claim

Questo capitolo non implica:

\begin{enumerate}
\def\labelenumi{\arabic{enumi}.}
\tightlist
\item
  che ogni struttura avanzata classica debba avere un corrispettivo ordinativo utile;
\item
  che la complessità strutturale sia di per sé un valore scientifico;
\item
  che i criteri qui proposti sostituiscano validazione empirica indipendente.
\end{enumerate}

\begin{center}\rule{0.5\linewidth}{0.5pt}\end{center}

\section{8. Claim Ledger (S0-S3)}

{\def\LTcaptype{none} % do not increment counter
\begin{longtable}[]{@{}lllll@{}}
\toprule\noalign{}
ID & Claim & Tipo & Evidenza & Stato \\
\midrule\noalign{}
\endhead
\bottomrule\noalign{}
\endlastfoot
C7.1 & Le strutture monoidali richiedono simmetria condizionata in domini asimmetrici & formale & S2 & in\_review \\
C7.2 & Le fibrations richiedono coerenza multilivello esplicita & formale-operativo & S1 & validated \\
C7.3 & La validità logica interna di un topos non implica validità ordinativa esterna & metateorico & S1 & validated \\
C7.4 & Le strutture avanzate possono mostrare deriva temporale pur restando formalmente lecite & formale-operativo & S2 & in\_proof \\
C7.5 & \texttt{Cost\_str} è utile come indicatore di controllo della complessità strutturale & metodologico & S1 & in\_review \\
C7.6 & Il criterio minimo di accettazione avanzata unifica monoidali/enriched/fibrations/topos & metodologico & S2 & in\_proof \\
C7.7 & Il capitolo chiude l'espansione classica \texttt{v1.3} in forma coerente con i blocchi precedenti & programmatico & S1 & in\_review \\
\end{longtable}
}

\begin{center}\rule{0.5\linewidth}{0.5pt}\end{center}

\section{9. Bridge al Capitolo 12}

Con questo capitolo si chiude la milestone \texttt{v1.3}:

\begin{enumerate}
\def\labelenumi{\arabic{enumi}.}
\tightlist
\item
  Capitolo 4: mattoni locali;
\item
  Capitolo 5: costrutti universali;
\item
  Capitolo 6: livello trasformazionale;
\item
  Capitolo 7: strutture avanzate e dinamiche categoriali.
\end{enumerate}

Il prossimo passaggio (\texttt{v1.4}) non aggiunge solo nuova teoria, ma standardizza il metodo:

\begin{enumerate}
\def\labelenumi{\arabic{enumi}.}
\tightlist
\item
  Capitolo 12 - protocollo scientifico OCT;
\item
  Capitolo 13 - benchmark e replicabilità;
\item
  Capitolo 14 - governance epistemica e quality gates.
\end{enumerate}

In sintesi: i capitoli 4-7 hanno esteso il linguaggio; i capitoli 12-14 ne definiranno la pratica scientifica condivisibile.

\begin{center}\rule{0.5\linewidth}{0.5pt}\end{center}

\section{Riferimenti}

Riferimenti di framework interno:

\begin{enumerate}
\def\labelenumi{\arabic{enumi}.}
\tightlist
\item
  Ghioni, F. (2026). \texttt{TE\_CORE\_v5.1.md}.
\item
  Ghioni, F. (2026). \texttt{Ordinative\_Set\_Theory\_OST\_A\_Concise\_Guide\_For\_AI\_v2\_1.md}.
\item
  \texttt{OCT\_FOUNDATIONAL\_BOOK\_CHAPTER\_06\_v1\_0.md}.
\item
  \texttt{OCT\_FOUNDATIONAL\_BOOK\_CHAPTER\_08\_v1\_0.md}.
\item
  \texttt{OCT\_FOUNDATIONAL\_BOOK\_CHAPTER\_11\_v1\_0.md}.
\end{enumerate}

Riferimenti primari:

\begin{enumerate}
\def\labelenumi{\arabic{enumi}.}
\tightlist
\item
  Mac Lane, S. (1998). \emph{Categories for the Working Mathematician}.
\item
  Riehl, E. (2016). \emph{Category Theory in Context}.
\item
  Spivak, D. (2014). \emph{Category Theory for the Sciences}.
\item
  Borceux, F. (1994). \emph{Handbook of Categorical Algebra}.
\end{enumerate}
